\section{Results}\label{sec:results}
\paragraph{Total error across all tasks.}
Figure~\ref{fig:holisticraw} ranks all 14 evaluated optimizer pipelines by summed raw extrapolation RMSE. Each task contributes its mean error across sizes and starts. Iterative WSP has the smallest sum in both information conditions, approximately $2.03\times10^9$. The examples-only one-shot sum is $1.06\times10^{12}$. More than 99.8\% of that one-shot total comes from Fibonacci, while positive self-power contributes 78.1\% of the iterative total. The aggregate therefore emphasizes errors on targets with large values. Figure~\ref{fig:holisticfull} shows the full range including divergent recurrent fits.

\paragraph{Central and extrapolation behavior.}
The median scaled error answers a different aggregation question. With examples alone, iteration reduces the held-out central or early median from 0.0214 to 0.000765, while the extrapolation median rises from 0.057 to 0.235. The scaled error divides each output residual by the larger of one and the target magnitude before computing RMSE. Figure~\ref{fig:priorscaledmedian} retains that view in the appendix. Table~\ref{tab:paired_methods} gives the eight primary endpoints.
\paragraph{What changes during revision?}
Figure~\ref{fig:wsp-learning} separates training fit, development feedback and untouched extrapolation across three revision opportunities. Development loss cannot rise under the selection rule, but neither training nor extrapolation error has that guarantee. With formulas and examples, median scaled outer RMSE falls from 0.0731 to 0.00590 after the first opportunity, then rises to 0.227 after the second. The final table alone hides this transient improvement. Test scores are retrospective and never select weights; stopped runs retain their last state. Appendix~\ref{app:learning-curves} separates proposed and retained weights and gives raw errors by task.
\input{learning_main}
\input{iterative_paired_table}
\paragraph{Does feedback improve extrapolation?}
With examples only, iterative WSP returned lower outer RMSE than one-shot WSP in \iterBetterE{} of 87 settings, the same RMSE in \iterTieE{}, and higher RMSE in \iterWorseE{}. With the formula supplied, the counts were \iterBetterF{}, \iterTieF{} and \iterWorseF{}. The number of settings exact on every outer reference changed from \oldPerfectE{} to \iterPerfectE{} with examples only and from \oldPerfectF{} to \iterPerfectF{} with the formula. In both conditions, the selected development loss fell strictly in 52 settings and tied in 35. In the examples-only condition the loss on the original training examples rose in 46, fell in six and tied in 35. The prompt requests exact unseen-domain outputs among similarly fitting solutions. The revision selector gives priority to central development loss. Figure~\ref{fig:prediction-curves} shows the computed functions for three fixed scalar cases. Across all six width-16 square networks, iteration moves every change of slope into $[-3,3]$, leaving affine predictions on the outer intervals. At input 10, one examples-only run changes its prediction from approximately 100 to 46.7. Appendix~\ref{app:prompts} reproduces the instructions, and Appendix~\ref{app:feedback} gives the coefficient-level diagnosis.

\paragraph{How does WSP compare with backpropagation-based updates?}
One-shot examples-only WSP gave lower outer RMSE than plain gradient descent in 75 of 87 settings and lower RMSE than gradient descent followed by L-BFGS in 75. For iterative examples-only WSP these counts were \vsGDE{} and \vsGLE{}, respectively. Neither training-only gradient procedure returned an entirely exact outer set. Refinement nevertheless mattered: at width two on subtraction, median RMSE fell from 0.831 after plain gradient descent to $1.44\times10^{-13}$ after L-BFGS. The two examples-only WSP procedures both matched all 2,048 inputs at each start. The original examples-only reply had lower outer RMSE than Adam followed by full BFGS in 73 settings and than trust-region least squares in 76. Its training loss was lower than Adam followed by BFGS in 79 settings. The full-information gradient control in Appendix~\ref{app:feedback} checks what changes if the numerical solvers also have access to the development labels.

\input{prediction_curves}

\paragraph{How large were the errors?}
Table~\ref{tab:focusouter} in the appendix gives the size of the errors and the exact-match counts at a specified network size for every task. Figures~\ref{fig:taskbars1}--\ref{fig:taskbars5} plot those same values for all eight methods. We chose sizes from the architecture and representability arguments: the largest tested scalar width, the two-unit arithmetic construction, and the smallest known sufficient real-arithmetic size for the sequence and sorting tasks. Table~\ref{tab:sizes_outer} gives all 29 pairs on the outer test. Tables~\ref{tab:sizes_dev} and~\ref{tab:sizes_train} report the central or early test and the training examples. No test score was used to pick an optimizer checkpoint.

\paragraph{Did the optimizer reach the global training minimum?}
For the four scalar tasks with one hidden unit we can compute the global optimum on the supplied labels in real arithmetic. After sorting the training inputs there are finitely many intervals for the change of slope. For each interval we fit the two affine pieces subject to continuity, also checking the boundaries, and take the smallest feasible loss. Appendix~\ref{app:optimizer} gives the rational arithmetic procedure for these same training data. This yields an independent benchmark in 12 settings. At an absolute tolerance of $10^{-10}$ above the certified optimum, one-shot examples-only WSP reached it in 1 of 12 settings, one-shot formula-and-examples WSP in 2, Adam followed by BFGS in 1, and trust-region least squares from Xavier in 5. Iterative WSP reached it in \certIterE{} and \certIterF{} settings in the two information conditions, respectively; gradient descent reached it in \certGD{} and gradient descent followed by L-BFGS in \certGL{}. This reverses the ordering of examples-only WSP and least squares seen in the outer-error comparisons.

\paragraph{What if least squares starts from the WSP weights?}
We took each of the 87 examples-only WSP arrays and applied the same trust-region solver and training-loss stopping rules used from Xavier. On the 12 certified fits this initialization brought the solver within tolerance of the global minimum in 12 of 12. Across the full benchmark the refinement lowered training loss by more than $10^{-10}$ in 22 settings. Among those 22, mean squared scaled outer error (with the same factor of one half as training) increased by more than $10^{-10}$ in 13 and decreased by more than $10^{-10}$ in nine. The 13 deteriorations also increased raw outer RMSE. The number of settings with every outer prediction exact stayed at 30. This two-stage experiment gives the refiner an additional model call through its starting weights. Zero extra numeric solver steps were available to the single-response WSP methods.

